\section{Examples and a collapsing bounded face}\label{sec:examples}

The following examples serve different purposes. Scalar sums give
explicit norm paths and a visible change of support topology. The
subsequent triples lie outside that simultaneous-apartment setting;
their paths are supplied by the theorem, without an asserted explicit
formula for their supports.

\subsection{The scalar junction map}
In the original coordinate put
\[
 \lambda(x,t)=\frac{1}{\sqrt{1+t^2}}
                   \arctan\frac{x}{\sqrt{1+t^2}},\qquad
 A_\mu(t)=\int_0^t\lambda(\mu s,s)\,\dd s.
\]
Thus $A_{\pm a}(t)=\pm F(\arsinh t)$ for $t\ge0$.
A scalar Y meeting at $(p,\tau)$ has incoming position $p$
and outgoing positions $p-a(t-\tau)$ and $p+a(t-\tau)$.
Given incoming logarithmic norm $b_0$, set
$k=b_0+\int_{-\infty}^\tau\lambda(p,s)\,\dd s$.
Its common incoming weight is
$b_0+\int_{-\infty}^t\lambda(p,s)\,\dd s$ up to $\tau$,
and its outgoing weights are
\[
 w_v(t)=k+\int_\tau^t\lambda(p+\epsilon_va(s-\tau),s)\,\dd s.
\]
These weights agree to first order at $\tau$, and their gap is
positive for $t>\tau$. Their normalized outgoing constants are
$b_0+d_A$ and $b_0+d_C$, where
\begin{equation}\label{eq:scalarboundarymap}
 d_v=\int_{-\infty}^\tau\lambda(p,s)\,\dd s
  +\int_\tau^\infty
       \bigl[\lambda(p+\epsilon_va(s-\tau),s)
                    -\lambda(\epsilon_va s,s)\bigr]\,\dd s
  -A_{\epsilon_va}(\tau).
\end{equation}
The integrands at the ends are $O(|s|^{-2})$.

\begin{proposition}\label{prop:scalarjunction}
The map $(p,\tau)\mapsto(d_A,d_C)$ is a smooth global
diffeomorphism of $\RR^2$. Thus each scalar full boundary triple
has one scalar Y, which is the path of Theorem~\ref{thm:main}.
\end{proposition}

\begin{proof}
Write $f(x,s)=(1+x^2+s^2)^{-1}=\partial_x\lambda(x,s)$,
and set
\[
 P=\int_{-\infty}^\tau f(p,s)\,\dd s,\quad
 L=\int_\tau^\infty f(p-a(s-\tau),s)\,\dd s,\quad
 R_+=\int_\tau^\infty f(p+a(s-\tau),s)\,\dd s.
\]
Differentiating the convergent integrals gives
\[
 d_{A,p}=P+L,\quad d_{A,\tau}=aL,\qquad
 d_{C,p}=P+R_+,\quad d_{C,\tau}=-aR_+.
\]
Their Jacobian is $-a\{P(L+R_+)+2LR_+\}<0$.
For a global argument use $y=a\tau$ and output coordinates
$U=d_A+d_C$, $V=d_A-d_C$. The derivative matrix is
\[
 H=\begin{pmatrix}2P+L+R_+&L-R_+\\ L-R_+&L+R_+\end{pmatrix},
 \qquad
 z^THz=2Pz_1^2+L(z_1+z_2)^2+R_+(z_1-z_2)^2.
\]
It is symmetric positive definite. Hence $(U,V)=\nabla G$
for a smooth strictly convex function $G$ on $\RR^2$.
Put $\rho=\sqrt{1+p^2+\tau^2}$. In each integral retain a
segment of length $\rho$ adjacent to $\tau$. On those segments
$1+x^2+s^2\le C\rho^2$, with a constant depending only on
$a$. Thus $P,L,R_+\ge c/\rho$, and
\[
 H\ge\frac{c'}{\sqrt{1+p^2+y^2}}I.
\]
At the origin $(U,V)=0$. Along every unit ray this implies
$e\cdot\nabla G(re)\ge c'\arsinh r$. Integration makes
$G(q)-w\cdot q$ coercive for every $w\in\RR^2$.
Its unique minimum solves $\nabla G(q)=w$. The inverse
function theorem gives a smooth inverse everywhere.
The scalar Y is a finite regular matched path with the prescribed
ends, so Theorem~\ref{thm:wholelineconverse} identifies it
with the constructed minimum.
\end{proof}

\subsection{An explicit degree-three bounded face}
Fix $p\ge0$ and $(r_1,r_2,r_3)=(-p,0,p)$. In the standard
basis of $K^3$ define
\begin{equation}\label{eq:diamondweights}
 w_{vi}(t;p)=
 \begin{cases}
 \displaystyle\int_{-\infty}^t\lambda(r_i,s)\,\dd s,&t\le0,\\[4pt]
 \displaystyle\int_{-\infty}^0\lambda(r_i,s)\,\dd s
       +\int_0^t\lambda(r_i+\epsilon_va s,s)\,\dd s,&t\ge0.
 \end{cases}
\end{equation}
Let $h_v(t;p)$ be the corresponding diagonal norms.
Their full incoming norm is the standard maximum norm $h_0$.
Their outgoing normalized norms are diagonal in that same basis,
with weights
\begin{equation}\label{eq:diamondboundary}
 b_{vi}(p)=\int_{-\infty}^0\lambda(r_i,s)\,\dd s
   +\int_0^\infty
      \bigl[\lambda(r_i+\epsilon_va s,s)
                     -\lambda(\epsilon_va s,s)\bigr]\,\dd s.
\end{equation}
These formulas specify the norms on every vector, including their
scales. Each coordinate is the scalar Y at $(r_i,0)$.
The pair is equal for $t\le0$ and ordered for $t\ge0$, since
\[
 w_{Ci}(t;p)-w_{Ai}(t;p)
 =\int_0^t[\lambda(r_i+as,s)-\lambda(r_i-as,s)]\,\dd s>0
 \quad(t>0).
\]
It has first-order matching at zero and smooth weights elsewhere.
Theorem~\ref{thm:wholelineconverse} identifies it with the
unique matched path for \eqref{eq:diamondboundary}.

For $p>0$ the Y vertices and crossing vertices are
\begin{equation}\label{eq:diamondvertices}
 \begin{gathered}
 Y_1=(-p,0),\quad Y_2=(0,0),\quad Y_3=(p,0),\\
 P_{12}=(-p/2,p/(2a)),\quad
 P_{23}=(p/2,p/(2a)),\quad P_{13}=(0,p/a).
 \end{gathered}
\end{equation}
Indeed $C_i:x=r_i+at$ meets $A_j:x=r_j-at$ for $t>0$
exactly when $i<j$. There are six vertices, six finite edges
and nine exterior rays, with every open edge of multiplicity
one. The unique bounded face has boundary
\[
 Y_2\longrightarrow P_{12}\longrightarrow P_{13}
          \longrightarrow P_{23}\longrightarrow Y_2.
\]
The two other finite edges connect $Y_1$ to $P_{12}$ and
$Y_3$ to $P_{23}$. The finite core is connected with one cycle;
there are no further crossings, so exactly one face is bounded.
Its Euclidean area in the $(x,t)$ plane is $p^2/(2a)$.
This is not its area for the form $\omega$.

The vertex object at $Y_i$ is $B$ on the $i$th coordinate.
It has incoming factor $B$ and outgoing filtration
$0\subset A\subset B$, with quotient $C$.
At $P_{ij}$ the object is $A_j\oplus C_i$ with zero arrow,
since $i\ne j$. Both filtrations have factors $A_j,C_i$
in that order. Thus a crossing does not create a $B$ strand.
These labels and filtrations give an explicit instance of
Proposition~\ref{prop:vertexfiltrations}.

\begin{figure}[t]
\centering
\begin{tikzpicture}[x=2cm,y=2cm,line width=.7pt,
  every node/.style={font=\small},>=Stealth]
  \begin{scope}
    \fill[blue!8] (0,0)--(-.5,.5)--(0,1)--(.5,.5)--cycle;
    \foreach \r in {-1,0,1} {
      \draw (\r,-.6)--(\r,0);
      \draw (\r,0)--({\r-1.45},1.45);
      \draw (\r,0)--({\r+1.45},1.45);
    }
    \foreach \x/\y in {-1/0,0/0,1/0,-.5/.5,.5/.5,0/1}
      \fill (\x,\y) circle (1.4pt);
    \node[below right] at (-1,0) {$Y_1$};
    \node[below right] at (0,0) {$Y_2$};
    \node[below left] at (1,0) {$Y_3$};
    \node[left] at (-.5,.5) {$P_{12}$};
    \node[right] at (.5,.5) {$P_{23}$};
    \node[above] at (0,1) {$P_{13}$};
    \node at (-1.8,1.35) {$A$};
    \node at (1.8,1.35) {$C$};
    \node[left] at (-1,-.45) {$B$};
    \node at (0,-.9) {$p>0$};
  \end{scope}
  \begin{scope}[xshift=7.3cm]
    \draw (0,-.6)--(0,0)--(-1.05,1.05);
    \draw (0,0)--(1.05,1.05);
    \fill (0,0) circle (1.6pt);
    \node[left] at (0,-.4) {$B^3$};
    \node[left] at (-.65,.65) {$A^3$};
    \node[right] at (.65,.65) {$C^3$};
    \node at (0,-.9) {$p=0$};
  \end{scope}
\end{tikzpicture}

\caption{The support of \eqref{eq:diamondweights} for $p>0$
(left) and $p=0$ (right). The left panel uses vertical coordinate
$a t/p$ and horizontal coordinate $x/p$. Its finite core has six
vertices and one bounded diamond. At $p=0$ the entire finite core collapses to
one Y with edge multiplicity three. The four-valent crossings have
vertex objects $A\oplus C$.}
\label{fig:diamondcollapse}
\end{figure}

\begin{proposition}\label{prop:collapsecontinuity}
For all $p,p'\ge0$,
\begin{equation}\label{eq:collapsecontinuity}
 \sup_{t\in\RR}\max_{v=A,C}
 d_2(h_v(t;p),h_v(t;p'))\le\pi\sqrt2\,|p-p'|.
\end{equation}
The same bound holds for each outgoing full boundary norm.
At $p=0$ all three boundary norms are $h_0$, and the support
is the single Y with labels $B^3,A^3,C^3$.
\end{proposition}

\begin{proof}
Since $0<\partial_x\lambda(x,s)\le(1+s^2)^{-1}$,
each integral in \eqref{eq:diamondweights} gives
\[
 |w_{vi}(t;p)-w_{vi}(t;p')|
       \le\pi|r_i(p)-r_i(p')|
\]
uniformly in $t$. The same basis diagonalizes both norms, so
summing the squares of the three weight differences proves
\eqref{eq:collapsecontinuity}. Multiplication by the same
reference scalar is an isometry; taking normalized limits
gives the outgoing bound. At $p=0$, the incoming weights are
zero and the outgoing weights are $A_{\epsilon_va}(t)$.
Normalization removes the latter exactly. This proves the
last assertion and displays the collapse explicitly.
\end{proof}

The estimate proves continuity of the full norm paths through
this particular topology change. It does not establish uniform
continuity for arbitrary boundary triples. Although this example
has a bounded face, its norms are a direct sum of scalar norms;
bounded faces do not require irreducible boundary data.

\subsection{A triple with no common orthogonal splitting}
Fix $u>0$. Give the following three bases the weights $(-2u,-u,0)$:
\begin{equation}\label{eq:irreducibletriple}
 \begin{split}
 b^{(1)}&=(e_1,e_2,e_3),\\
 b^{(2)}&=(e_3,e_2,e_1),\\
 b^{(3)}&=(e_1+e_2+e_3,\ e_1+2e_2+3e_3,\ e_3).
 \end{split}
\end{equation}
Their basis matrices are invertible over $\CC$. Denote the
resulting norms by $h_1,h_2,h_3$ and the standard norm by $h_0$.
They are valid full boundary data $(h_B,h_A^+,h_C^+)$.

We verify that these norms have neither a common apartment nor
a nontrivial common orthogonal direct-sum decomposition.
Each is at most $h_0$. In its residue space the vectors on
which $h_i<h_0$ span the plane generated by the first two
vectors of $b^{(i)}$. These planes are
\[
 P_1=\langle e_1,e_2\rangle,\quad
 P_2=\langle e_2,e_3\rangle,\quad
 P_3=\{x_1-2x_2+x_3=0\}.
\]
Their common intersection is zero. Rescale any nonzero vector
to have $h_0$-norm one, using the full real value group.
Its nonzero residue misses one of the planes, so some $h_i$
has value one on it. Hence $\max_i h_i=h_0$ pointwise.

If the three norms had a common apartment, their maximum
would belong to it too. Reducing a common $h_0$-orthonormal
basis would split all three complete residue flags. The
first two flags are opposite, with canonical intersection
lines $\CC e_1,\CC e_2,\CC e_3$. A basis splitting both
must consist of these lines. The first line of the third
flag is $\CC(e_1+e_2+e_3)$, a contradiction.

More generally, a common orthogonal decomposition would be
orthogonal for $h_0$ and would reduce to a decomposition of
$\CC^3$ splitting the three flags. Its projection preserves
the first two opposite flags, so it is a diagonal idempotent
in the coordinate basis. Preserving the third flag's first
line forces all three diagonal entries to agree. The
projection is therefore zero or identity, and one residue
summand, hence one original summand, is zero. This excludes
a nontrivial decomposition. Theorem~\ref{thm:main} applies
to these triples without reducing them to scalar paths.

\subsection{Full norm information beyond pairwise spectra}
Here is a valued family making that distinction numerical.
Set $u=1$, take $0<s<1$, and put $b=T^s$. Keep the first
two bases in \eqref{eq:irreducibletriple}, and replace the
third by the columns of
\[
 B_3(s)=\begin{pmatrix}1&1&0\\b&2&0\\1&3&1\end{pmatrix}.
\]
Its determinant $2-b$ is a unit. Give all three bases weights
$(-2,-1,0)$. Every pairwise logarithmic relative spectrum is
$(2,0,-2)$, independently of $s$. One can see this from the
exterior minors: after including the norm weights the minimum
valuations of the entries and of the two-by-two minors are
both $-2$, while the determinant valuation is zero. Unit
extremal minors attain these minima; positive-valuation
entries cannot decrease them.

For an explicit three-norm invariant, use the standard volume
form and let $L_i=\{v:h_i(v)\le1\}$. Define
\[
 D_{111}(h_1,h_2,h_3)
   =-\min\{\nu\det(v_1,v_2,v_3):v_i\in L_i\}.
\]
Each lattice has the finite orthonormal generators
$T^{-2}b^{(i)}_1,T^{-1}b^{(i)}_2,b^{(i)}_3$.
Multilinearity and the ultrametric inequality reduce the
minimum to the finitely many choices of one generator from
each lattice. Choosing the first generator in each gives
valuation $-6+s$, because the unscaled determinant is
$\pm b$. Every other choice has weight sum at least $-5$
and an integral determinant. Since $s<1$, the minimum is
$-6+s$, and therefore
\begin{equation}\label{eq:tripleinvariant}
 D_{111}=6-s.
\end{equation}
Thus the pairwise spectra discard information retained by
the full norm triple. The general theorem applies to every
member of this valued family. We do not identify its support
with the diamond above, or with the dual of a numerical hive
built from determinant invariants.


\subsection{Reproducibility}
The accompanying source includes the script
\texttt{check\_examples.py}. It enumerates the nonzero exterior minors of the valued family
and verifies their affine valuation inequalities on the whole
interval $0<s<1$. It also computes the triple determinant
minimum and the rational coordinates, incidence relations and
area of the normalized diamond. Figure~\ref{fig:diamondcollapse}
is drawn directly from those coordinates in the accompanying
TikZ source. The analytic convergence and continuity arguments
are the proofs above and do not depend on quadrature.
